\section{A uniform Boolean sampler}
\label{sec:sampler}

We next construct an averaging operator from random subset sums. The
exponential-potential argument develops the finite-group binary correction
of Li and Liu~\cite[Theorem~4.1]{LLC}; we prove the weighted estimate with
random shifts used below, including all required quantitative bounds. A second
moment estimate alone gives a good approximation at most base points. A
second group of random shifts upgrades this to a lower bound at \emph{every}
base point. This distinction will allow a later block to sample an earlier
one even when its base point depends on other lattice choices.

\begin{lemma}[Uniform Boolean sampling]
\label{lem:sampler}
There is an absolute constant $b_0$ with the following property. Let
$b\ge b_0$, let $H$ be a finite abelian group written additively, and suppose
that
\[
  \log |H|\le b^4,
  \qquad 0\le g\le \exp(b^{0.65}),
  \qquad \bar g:=\E_{x\in H}g(x)\in[1/2,2].
\]
Let $s$ be an integer satisfying $b^{0.69}\le s\le b$, and choose
$w_1,\ldots,w_s$ independently and uniformly in $H$. With
$c_*=(\log 2)/8$, the probability that
\begin{equation}
\label{eq:sampler-conclusion}
  2^{-s}\sum_{\eps\in\{0,1\}^s}
  g\left(x-\sum_{i=1}^s\eps_iw_i\right)
  \ge (1-3b^{-1/10})\bar g
  \qquad\text{for every }x\in H
\end{equation}
is at least
\begin{equation}
\label{eq:sampler-probability}
  1-\exp(-c_*b^{0.69})
    -b\exp\left(-\frac{c_*}{2}b^{0.59}\right).
\end{equation}
In particular, the failure probability is at most $b^{-1/2}$ after
increasing $b_0$.
\end{lemma}

\begin{proof}
Put $\delta=b^{-1/10}$, $k=\lfloor s/2\rfloor$, and $j_*=s-k$.
We use the first $k$ shifts to define
\[
 A(x)=2^{-k}\sum_{\eps\in\{0,1\}^k}
       g\left(x-\sum_{i=1}^k\eps_iw_i\right).
\]
Temporarily take $x$ independently uniform in $H$. For two distinct bit
vectors $\eps,\eps'$, their arguments in this average are independent
uniform points. To see this, choose an index at which the bits differ and
condition on all other shifts. The map from $x$ and that remaining shift
to the two arguments is a bijection of $H\times H$: subtracting the
arguments determines the shift with coefficient $1$ or $-1$, and then
determines $x$. This works in every finite abelian group, with no
restriction on its order. Consequently,
\begin{equation}
\label{eq:sampler-second-moment}
 \E_{w_1,\ldots,w_k,x}(A(x)-\bar g)^2
 =2^{-k}\left(\E_xg(x)^2-\bar g^2\right)
 \le 2^{-k}\exp(2b^{0.65}).
\end{equation}

For fixed first-stage shifts, let
\[
 \mathcal G=\{x\in H:A(x)\ge(1-\delta)\bar g\},
 \qquad \beta=1-\frac{|\mathcal G|}{|H|}.
\]
Chebyshev's inequality, $\bar g\ge1/2$, and
Equation~\eqref{eq:sampler-second-moment} give
\[
 \E_{w_1,\ldots,w_k}\beta
 \le 4b^{1/5}2^{-k}\exp(2b^{0.65})
 \le \exp(-2c_*s).
\]
The last inequality holds for all sufficiently large $b$, uniformly in
$s\ge b^{0.69}$. Indeed, its logarithm is at most
\[
 \log 8+\tfrac15\log b+2b^{0.65}-\tfrac12s\log 2,
\]
and the first three terms are at most $s\log 2/4$ once $b$ is large.
Markov's inequality therefore yields
\begin{equation}
\label{eq:sampler-small-bad-set}
 \Prob\{\beta>\exp(-c_*s)\}\le\exp(-c_*s).
\end{equation}
Fix any first-stage shifts for which
$\beta\le\exp(-c_*s)$.

For $0\le j\le j_*$, let
\[
 N_j(x)=
 \sum_{\eta\in\{0,1\}^j}
 \ind_{\mathcal G}\left(x-\sum_{i=1}^j\eta_iw_{k+i}\right).
\]
Here we count bit labels with multiplicity, even if their subset sums
coincide. Thus, for a fresh shift $w=w_{k+j+1}$,
\[
 N_{j+1}(x)=N_j(x)+N_j(x-w).
\]
Set $\theta=c_*s/2$ and introduce the potential
\[
 \Phi_j=\E_{x\in H}\exp(-\theta N_j(x)).
\]
Since $N_0=\ind_{\mathcal G}$, the bound on $\beta$ gives
\[
 \Phi_0=\beta+(1-\beta)e^{-\theta}
 \le e^{-2\theta}+e^{-\theta},
 \qquad \log\Phi_0\le-\theta+1.
\]
Most importantly, conditional on every shift chosen so far,
\begin{equation}
\label{eq:sampler-potential-square}
 \E_w\Phi_{j+1}
 =\E_{x,w}e^{-\theta N_j(x)}e^{-\theta N_j(x-w)}
 =\Phi_j^2.
\end{equation}
The last equality uses the bijection $(x,w)\mapsto(x,x-w)$ of
$H\times H$. It makes no independence assumption about the subset-sum
labels themselves.

Conditional Markov's inequality applied to
Equation~\eqref{eq:sampler-potential-square} shows that
\[
 \Prob\left\{
    \log\Phi_{j+1}>2\log\Phi_j+\theta\delta
    \,\middle|\,\text{previous shifts}
 \right\}
 \le e^{-\theta\delta}.
\]
There are $j_*\le b$ steps. Except on an event of conditional probability
at most $b e^{-\theta\delta}$, the recursion holds at every step and gives
\begin{equation}
\label{eq:sampler-potential-iterate}
 \log\Phi_j
 \le 2^j(-\theta+1)+(2^j-1)\theta\delta
 \qquad(0\le j\le j_*).
\end{equation}
For each $x$, its single summand in the potential implies
\[
 \Phi_j\ge |H|^{-1}e^{-\theta N_j(x)}.
\]
Combining this with Equation~\eqref{eq:sampler-potential-iterate} yields
\begin{equation}
\label{eq:sampler-uniform-good-fraction}
 \frac{N_{j_*}(x)}{2^{j_*}}
 \ge 1-\delta-\frac1\theta
             -\frac{\log|H|}{\theta 2^{j_*}}
 \ge 1-2\delta
 \qquad(x\in H).
\end{equation}
For the last inequality, use
$\theta\ge(c_*/2)b^{0.69}$, $j_*\ge b^{0.69}/2$, and
$\log|H|\le b^4$. Uniformly under these bounds,
$1/\theta\le\delta/2$ and
$\log|H|/(\theta2^{j_*})\le\delta/2$ once $b_0$ is sufficiently large.

Finally, the full average factors as
\[
 2^{-s}\sum_{\eps\in\{0,1\}^s}
 g\left(x-\sum_{i=1}^s\eps_iw_i\right)
 =2^{-j_*}\sum_{\eta\in\{0,1\}^{j_*}}
 A\left(x-\sum_{i=1}^{j_*}\eta_iw_{k+i}\right).
\]
On labels counted by $N_{j_*}(x)$, the summand is at least
$(1-\delta)\bar g$; all other summands are nonnegative. Thus
Equation~\eqref{eq:sampler-uniform-good-fraction} bounds this expression
below by
\[
 (1-2\delta)(1-\delta)\bar g
 \ge(1-3\delta)\bar g,
\]
simultaneously for every $x$. By
Equation~\eqref{eq:sampler-small-bad-set}, the total failure probability
is at most
\[
 e^{-c_*s}+b e^{-\theta\delta}
 \le e^{-c_*b^{0.69}}
      +b e^{-(c_*/2)b^{0.59}},
\]
as asserted. All thresholds used above are absolute.
\end{proof}

\begin{remark}[Uniformity and offsets]
The first stage controls only the proportion of bad base points. A direct
union bound over $H$ would be insufficient under the allowed bound
$\log|H|\le b^4$. The second-stage potential instead pays
$\log|H|$ in Equation~\eqref{eq:sampler-uniform-good-fraction}, divided by
an exponentially large number of labels. Once its conclusion holds, the
base point can be replaced by any element of $H$, including an offset
determined by other random columns. In particular, suppose a subset is
chosen independently of the column values and its size lies in the range
of Lemma~\ref{lem:sampler}. One may condition on that subset, apply the
lemma to its columns, and then insert arbitrary offsets formed from the
complementary columns. This observation will justify the biased-bit
mixtures below.
\end{remark}
